\documentclass[11pt]{article}
\usepackage{amsmath,amssymb,amsthm}
\newtheorem{theorem}{Theorem}
\newtheorem{corollary}{Corollary}
\title{A Compact Shared-Residue Pool Cannot Supply Two Arbitrary Words}
\author{}
\date{}
\begin{document}
\maketitle

\paragraph{Status and scope.}
This is a complete written proof, not a Lean-checked result. It concerns
a square exterior Cauchy decoder with shared absolute residues, and
does not give a lower bound for unrestricted HeadComplexity. It extends
the explicit three-pole calculation in
\texttt{common-cauchy-amplitude-obstruction.tex}. The compactness
threshold below depends on the fixed number of poles and the fixed
query grids; no quantitative threshold uniform in their sizes is
claimed.

Fix an integer $q\geq3$ and two strictly increasing positive query
grids $u_1<\cdots<u_q$ and $v_1<\cdots<v_q$. For distinct poles
$0<p_1<\cdots<p_q$, write
\[
 P(t)=\prod_{h=1}^q(t-p_h),\qquad
 F(t)=\sum_{h=1}^q\frac{a_h}{t-p_h}=\frac{N(t)}{P(t)}.
 \tag{1}
\]
The numerator in (1) has degree at most $q-1$. A sign change of a
strict sign word means an adjacent pair of entries with opposite
signs.

\begin{theorem}[Proper flips lose exterior oscillations]
There is a $\delta>0$, depending only on $q,u,v$, with the following
property. Suppose $p_q<\delta$ and $F(u_i)$ has $q-1$ sign changes.
For any $\epsilon_h\in\{-1,1\}$ which are not all equal, put
\[
 F_\epsilon(t)=\sum_{h=1}^q\frac{\epsilon_h a_h}{t-p_h}.
 \tag{2}
\]
If all $F_\epsilon(v_i)$ are nonzero, their sign word has at most
$\lfloor q/2\rfloor-1$ sign changes. This conclusion is uniform
over the poles and the residues $a_h$.
\end{theorem}

\begin{proof}
Shrink $\delta$ so all poles are below both query grids. The
denominator $P$ is positive on both grids. Full alternation at $u$
forces $N$ to have exactly $q-1$ real roots, one in each open interval
$(u_i,u_{i+1})$, and to have degree $q-1$. Multiplying every residue
by the same nonzero scalar does not change either number of sign
changes, so normalize $N$ to be monic. Thus
\[
 N(t)=\prod_{i=1}^{q-1}(t-\rho_i),\qquad
 u_i<\rho_i<u_{i+1}.
 \tag{3}
\]
These normalized polynomials lie in a compact set of coefficient
vectors, and their constant coefficients are uniformly bounded away
from zero. Residue evaluation gives
\[
 a_h=\frac{N(p_h)}{P'(p_h)}.
 \tag{4}
\]

Suppose the theorem were false. There would be a sequence of such
configurations with $p_q\longrightarrow0$ and proper flips producing
at least $k=\lfloor q/2\rfloor$ sign changes at $v$. Pass to
subsequences so the flip vector is fixed, the set of $k$ adjacent
intervals with a sign change is fixed, and $N$ converges in
coefficients to a monic polynomial $N_0$ of the form (3) with closed
root intervals. In particular $N_0(0)\neq0$.

Let $R$ be the numerator of $F_\epsilon$. Then
\[
 R(p_h)=\epsilon_h N(p_h),\qquad
 P\mid(R^2-N^2),\qquad \deg R\leq q-1.
 \tag{5}
\]
The coefficient norm $\|R\|_\infty$ tends to infinity along this
subsequence. Otherwise a further bounded subsequence would converge
to a polynomial $R_*$. Evaluating at all $p_h\longrightarrow0$
would give $R_*(0)=\epsilon_h N_0(0)$ for every $h$, impossible
because both signs occur in $\epsilon$.

Normalize $\widetilde R=R/\|R\|_\infty$. Pass to a further
subsequence with coefficient limit $R_0$ of norm one. The remainders
in division by a monic polynomial are continuous functions of all
coefficients. Divide (5) by $\|R\|_\infty^2$ and let the poles
tend to zero. Since $P\longrightarrow t^q$ and
$N/\|R\|_\infty\longrightarrow0$, this gives
\[
 t^q\mid R_0(t)^2,
 \qquad
 t^{\lceil q/2\rceil}\mid R_0(t).
 \tag{6}
\]
Consequently the nonzero polynomial $R_0$ has at most
$q-1-\lceil q/2\rceil=\lfloor q/2\rfloor-1$ roots away from
zero, counted with multiplicity.

For completeness, the $k$ strict sign changes cannot disappear when
passing to the limit. Choose one root of $R$ in each of the fixed
$k$ open adjacent query intervals. Their monic product $J$ has degree
$k$ and divides $\widetilde R$. Pass to a subsequence on which all
chosen roots converge. Then $J$ converges to a monic polynomial
$J_0$ whose $k$ roots, counted with multiplicity, belong to
$[v_1,v_q]$. Monic division makes the quotient coefficients bounded
and continuous, so $J_0\mid R_0$. This argument retains multiplicity
when roots from neighboring intervals converge to the same shared
endpoint. Since $v_1>0$, it contradicts (6).
\end{proof}

\begin{corollary}[Two words on one grid already fail]
Fix $q\geq3$ and $0<u_1<\cdots<u_q$. For all sufficiently small
distinct positive poles $p_h$, there are two strict sign words on
this grid which cannot both be obtained from
\[
 \left(\sum_{h=1}^q\frac{\eta_h v_h}{u_i-p_h}\right)_{i=1}^q,
 \qquad \eta_h\in\{-1,1\},
 \tag{7}
\]
using any single positive amplitude vector $(v_h)$. The amplitude
vector is allowed to depend on the requested pair of words.
\end{corollary}

\begin{proof}
Request a fully alternating word and a word with exactly $q-2$
sign changes. If the first is represented, let $a_h$ be its signed
residues. Every representation of the second from the same amplitude
vector is a flip (2). A constant flip gives the first word or its
negative and therefore still has $q-1$ changes on this same grid.
A proper flip has at most $\lfloor q/2\rfloor-1$ changes by the
theorem. Both alternatives exclude the second word because
$q-2>\lfloor q/2\rfloor-1$ for $q\geq3$.
\end{proof}

\paragraph{Why the scope matters.}
The theorem's proper-flip bound permits a different second fixed
query grid. The two-word corollary uses the same grid: on arbitrary
different grids the unflipped numerator itself can have different
numbers of sign changes, so the corollary does not automatically
extend to that setting. Neither statement covers a rectangular pool
with more poles than queries, group-dependent absolute residues,
independently chosen numerator profiles, or noncompact pole
geometry. Even for square pools, the theorem alone supplies no
uniform compactness rate as $q$ grows. These are restrictions of a
specific proposed decoder, not obstructions to all admissible head
representations.
\end{document}
